% Created 2026-05-11 Mon 12:31
% Intended LaTeX compiler: pdflatex
\documentclass[11pt]{article}

\usepackage[utf8]{inputenc}
\usepackage[T1]{fontenc}
\usepackage{graphicx}
\usepackage{longtable}
\usepackage{wrapfig}
\usepackage{rotating}
\usepackage[normalem]{ulem}
\usepackage{amsmath}
\usepackage{amssymb}
\usepackage{capt-of}
\usepackage{hyperref}
\author{flavio}
\date{\today}
\title{Array}
\hypersetup{
 pdfauthor={flavio},
 pdftitle={Array},
 pdfkeywords={},
 pdfsubject={},
 pdfcreator={},
 pdflang={English}}
\begin{document}

\maketitle
\tableofcontents

Un array in java rappresenta un gruppo di variabili (chiamate elementi) tutte dello stesso tipo.

Gli array sono oggetti quindi le variabili di tipo array contengono il riferimento all'array in memoria. Gli elementi di un array possono essere tipi primitivi o riferimenti a oggetti.

Dichiarazione di un array senza valori. Ogni elemento è inzializzato con il valore di default.

\begin{minted}[]{java}
int[] a; // dichiarazione
a = new int[10] // creazione senza valori

\end{minted}

Creazione con valori
\begin{minted}[]{java}
a = new int[] {-5, -2, 3, 0, 1, -6, 75, 32, 122, 4};
\end{minted}

codice tipo
int[] a = new int[10] \{1,2,3,4,5,6,7,8,9,10\}

System.out.println(a[5])
Indice sempre positivo, compreso tra 0 e la dimensione dell'array-1
\subsection{Lughezza array}
\label{sec:org2cd5b67}
per la lughezza array \texttt{nomeArray.length} uh oh: campo pubblico? protected? coding horror.
Esempio di uso:
\begin{minted}[]{java}
public int[] getArrayFattorial(final int n)
    {
        int[] fatt = new int[n];
        fatt[0] = 1;
        for (int k=1; k < n; k++) fatt[k] = fatt[k-1] *k;
        return fatt;
    }
\end{minted}
\subsection{Le dimensioni di un Array non sono modificabili}
\label{sec:orgc43e93c}
Una volta dichiarate non crescono dinamicamente, tipo vettore in C++. Bisogna utilizzare il metodo \textbf{CopyOf}. Se le dimensioni dell'array sono maggiori di quelle iniziali allora setta al valore di default per il tipo.
\begin{minted}[]{java}
int[] array = {1,5,8,2,3,4}
    System.out.println(Arrays.toString(array));
array = Arrays.copyOf(array,5); //restringe
System.out.println(Arrays.toString(array));
array = Arrays.copyOf(array,8); //restringe
System.out.println(Arrays.toString(array));
\end{minted}
\section{Array a 2 dimensioni}
\label{sec:orgc937843}
Possiamo specificare un array a 2 dimensioni semplicemente specificando due coppie di parentesi quadre []:
String[][] matrice = new String[RIGHE][COLONNE]
L'accesso avviene specificando le due dimensioni dell'array:
matrice[y][x].
\end{document}
